\PassOptionsToPackage{table}{xcolor}
\documentclass[10pt,aspectratio=169]{beamer}
\newcommand{\LectureNumber}{24}
\input{course-theme.tex}
\title{Exception handling}
\subtitle{CSC103 Programming Fundamentals / Lecture 24}
\author{Dr. Muhammad Farhan}
\institute{Associate Professor of Computer Science\\Department of Computer Science\\COMSATS University Islamabad\\Sahiwal Campus}
\date{Fall 2026}
\begin{document}
\begin{frame}\titlepage\end{frame}

\begin{frame}[fragile,t]{The problem for this lecture}
\begin{columns}[T,onlytextwidth]\begin{column}{0.60\textwidth}
Define validateMark(int marks) that throws IllegalArgumentException outside 0..100. Test -1, 0, 100 and 101.
\end{column}\begin{column}{0.35\textwidth}\small
Write a validator that throws when marks are outside 0..100.
\end{column}\end{columns}
\note{Introduce the target problem before explaining the tools. Revisit it in the guided solution later.\par Syllabus reading: Deitel Ch 11. CLO-3.}
\end{frame}

\begin{frame}[fragile,t]{Learning objectives}
\begin{columns}[T,onlytextwidth]\begin{column}{0.60\textwidth}
\begin{itemize}
\item Explain propagation and the exception hierarchy
\item Use try, throw and catch
\item Distinguish checked and unchecked exceptions
\end{itemize}
\end{column}\begin{column}{0.35\textwidth}\small
CLO-3

Exceptional control flow\par Exception hierarchy\par Checked and unchecked exceptions\par Throwing and catching
\end{column}\end{columns}
\note{Explain how each objective supports the opening problem. The final questions check these objectives.\par Syllabus reading: Deitel Ch 11. CLO-3.}
\end{frame}

\begin{frame}[fragile,t]{Exceptional control flow}
\begin{itemize}
\item An exception interrupts the normal flow at the failure point.
\item Java searches for an applicable handler.
\item Statements after the failure in the same try body are skipped.
\end{itemize}
\vspace{1mm}\begin{block}{Example}
\begin{lstlisting}
try {
  int n = Integer.parseInt("cat");
  System.out.println(n);
} catch (NumberFormatException ex) {
  System.out.println("Invalid integer");
}
\end{lstlisting}
\end{block}
\note{The println(n) statement does not execute. A catch handler can recover, report or translate an error depending on the method contract.\par Syllabus reading: Deitel Ch 11. CLO-3.}
\end{frame}

\begin{frame}[fragile,t]{Exception hierarchy}
\begin{itemize}
\item Throwable is the common base for Error and Exception.
\item RuntimeException is a subclass of Exception.
\item Most application recovery targets specific Exception subclasses.
\end{itemize}
\vspace{1mm}\begin{block}{Example}
\begin{lstlisting}
Throwable
  Error
  Exception
    RuntimeException
\end{lstlisting}
\end{block}
\note{Checked exceptions exclude RuntimeException subclasses and Error subclasses. Do not normally catch Error to continue as if an ordinary input mistake occurred.\par Syllabus reading: Deitel Ch 11. CLO-3.}
\end{frame}

\begin{frame}[fragile,t]{Checked and unchecked exceptions}
\begin{itemize}
\item Checked exceptions must be caught or declared with throws.
\item Unchecked exceptions do not have that compiler requirement.
\item The distinction is about checking, not whether an event occurs at runtime.
\end{itemize}
\vspace{1mm}\begin{block}{Example}
\begin{lstlisting}
Checked: IOException
Unchecked: NumberFormatException
Unchecked: IllegalArgumentException
\end{lstlisting}
\end{block}
\note{Both kinds can occur during execution. A method should not suppress an exception simply to satisfy the compiler.\par Syllabus reading: Deitel Ch 11. CLO-3.}
\end{frame}

\begin{frame}[fragile,t]{Throwing and catching}
\begin{itemize}
\item throw raises a particular exception object.
\item throws appears in a method declaration.
\item Place specific catches before broader compatible catches.
\end{itemize}
\vspace{1mm}\begin{block}{Example}
\begin{lstlisting}
if (marks < 0 || marks > 100) {
  throw new IllegalArgumentException(
      "marks outside 0..100");
}
\end{lstlisting}
\end{block}
\note{A catch for Exception before a catch for NumberFormatException makes the latter unreachable. Introduce the validation rule as part of the method contract.\par Syllabus reading: Deitel Ch 11. CLO-3.}
\end{frame}

\begin{frame}[fragile,t]{Validation and parsing solve different problems}
\begin{itemize}
\item Parsing checks whether text can represent a number of the requested kind.
\item Validation checks whether the parsed number belongs to the application's allowed domain.
\item The value 101 parses as int but violates a marks range of 0..100.
\end{itemize}
\vspace{1mm}\begin{block}{Example}
\begin{lstlisting}
"cat": NumberFormatException while parsing
"101": parses, then fails range validation
"100": passes both checks
\end{lstlisting}
\end{block}
\note{Explain the mechanism using the displayed example. Parsing checks whether text can represent a number of the requested kind. Validation checks whether the parsed number belongs to the application's allowed domain. The value 101 parses as int but violates a marks range of 0..100.\par Syllabus reading: Deitel Ch 11. CLO-3.}
\end{frame}

\begin{frame}[fragile,t]{Exception handling changes the path}
\begin{itemize}
\item throw immediately interrupts the ordinary path in the current block.
\item A compatible catch can report the failure and allow later work to continue.
\item A loop with a try-catch inside its body can examine later inputs after one failure.
\end{itemize}
\vspace{1mm}\begin{block}{Example}
\begin{lstlisting}
For inputs -1,0,100,101:
Reject the first, accept the next two, reject the last.
\end{lstlisting}
\end{block}
\note{Explain the mechanism using the displayed example. throw immediately interrupts the ordinary path in the current block. A compatible catch can report the failure and allow later work to continue. A loop with a try-catch inside its body can examine later inputs after one failure.\par Syllabus reading: Deitel Ch 11. CLO-3.}
\end{frame}


\begin{frame}[fragile,t]{An exception transfers control to a handler}
\begin{center}\begin{tikzpicture}
\node[decision] (a) at (0,0) {parse succeeds?};
\node[alt] (b) at (-3.6,-1.65) {Use integer};
\node[hot] (c) at (3.6,-1.65) {Matching catch};
\node[box] (d) at (0,-3.25) {Continue after try/catch};
\draw[arr] (a.west) -| node[lab,pos=.25]{true} (b.north);\draw[arr] (a.east) -| node[lab,pos=.25]{false} (c.north);\draw[arr] (b.south) |- (d.west);\draw[arr] (c.south) |- (d.east);
\end{tikzpicture}\end{center}
\begin{block}{What to notice}When parsing fails, the remaining statements in that try block are skipped.\end{block}
\note{Trace each labelled connection. When parsing fails, the remaining statements in that try block are skipped.}
\end{frame}
\begin{frame}[fragile,t]{Key distinctions}
\small\renewcommand{\arraystretch}{1.5}
\rowcolors{2}{pale}{white}
\begin{tabularx}{\textwidth}{>{\raggedright\arraybackslash}X>{\raggedright\arraybackslash}X>{\raggedright\arraybackslash}X>{\raggedright\arraybackslash}X}
\rowcolor{navy}
\color{white}\bfseries Input & \color{white}\bfseries Parses as int? & \color{white}\bfseries Valid mark? & \color{white}\bfseries Outcome\\
"cat" & No & Not evaluated & Parsing failure\\
"-1" & Yes & No & Range failure\\
"0" & Yes & Yes & Accept\\
"100" & Yes & Yes & Accept\\
"101" & Yes & No & Range failure\\
\end{tabularx}
\vspace{3mm}\footnotesize These distinctions determine which operation or control structure fits the problem.
\note{\par Syllabus reading: Deitel Ch 11. CLO-3.}
\end{frame}

\begin{frame}[fragile,t]{First example: execution}
\begin{lstlisting}
try {
  int marks = Integer.parseInt("bad");
  System.out.println(marks);
} catch (NumberFormatException ex) {
  System.out.println("Invalid integer");
}
System.out.println("Finished");
\end{lstlisting}
\note{This is the main body from Java\_Examples/Lecture24.java. The source file includes the complete class. Predict the result before the next slide.\par Syllabus reading: Deitel Ch 11. CLO-3.}
\end{frame}

\begin{frame}[fragile,t]{First example: result}
\begin{columns}[T,onlytextwidth]\begin{column}{0.60\textwidth}
Invalid integer\par Finished\par The handler runs, then normal execution continues after the try-catch.
\end{column}\begin{column}{0.35\textwidth}\small
Checked and unchecked exceptions

Throwing and catching
\end{column}\end{columns}
\note{Expected output and interpretation: Invalid integer
Finished
The handler runs, then normal execution continues after the try-catch.\par Syllabus reading: Deitel Ch 11. CLO-3.}
\end{frame}

\begin{frame}[fragile,t]{Guided practice}
\begin{columns}[T,onlytextwidth]\begin{column}{0.60\textwidth}
Define validateMark(int marks) that throws IllegalArgumentException outside 0..100. Test -1, 0, 100 and 101.
\end{column}\begin{column}{0.35\textwidth}\small
Attempt the solution before continuing.

Use the definitions and examples from this lecture.
\end{column}\end{columns}
\note{Give students 8 minutes to attempt the problem. The following slides provide a complete solution. Original solution guidance: Throw when marks\textless{}0 || marks\textgreater{}100. Inputs -1 and 101 fail. Inputs 0 and 100 return normally.\par Syllabus reading: Deitel Ch 11. CLO-3.}
\end{frame}

\begin{frame}[fragile,t]{Solution strategy}
\begin{columns}[T,onlytextwidth]\begin{column}{0.60\textwidth}
\begin{itemize}
\item Write a validator that throws when marks are outside 0..100.
\item Call the validator for each boundary example inside a try block.
\item Print a success message only after validation returns normally.
\end{itemize}
\end{column}\begin{column}{0.35\textwidth}\small
The next program uses fixed inputs so its result can be reproduced.
\end{column}\end{columns}
\note{Discuss why each step is needed. State the input assumptions before executing the program.\par Syllabus reading: Deitel Ch 11. CLO-3.}
\end{frame}

\begin{frame}[fragile,t]{Complete solution (1/2)}
\begin{lstlisting}
public class Solution24 {
  public static void main(String[] args) throws Exception {
    int[] cases = {-1, 0, 100, 101};
    for (int mark : cases) {
      try {
        validateMark(mark);
        System.out.println(mark + " accepted");
      } catch (IllegalArgumentException ex) {
        System.out.println(mark + " rejected");
\end{lstlisting}
\note{Complete runnable file: Java\_Examples/Solution24.java. Inputs are fixed in the program.  \par Syllabus reading: Deitel Ch 11. CLO-3.}
\end{frame}

\begin{frame}[fragile,t]{Complete solution (2/2)}
\begin{lstlisting}
      }
    }
  }
  static void validateMark(int mark) {
    if (mark < 0 || mark > 100) {
      throw new IllegalArgumentException("marks outside 0..100");
    }
  }
}
\end{lstlisting}
\note{Complete runnable file: Java\_Examples/Solution24.java. Inputs are fixed in the program.  \par Syllabus reading: Deitel Ch 11. CLO-3.}
\end{frame}

\begin{frame}[fragile,t]{Execution trace}
\small\renewcommand{\arraystretch}{1.5}
\rowcolors{2}{pale}{white}
\begin{tabularx}{\textwidth}{>{\raggedright\arraybackslash}X>{\raggedright\arraybackslash}X>{\raggedright\arraybackslash}X>{\raggedright\arraybackslash}X}
\rowcolor{navy}
\color{white}\bfseries Mark & \color{white}\bfseries Invalid condition & \color{white}\bfseries Control path & \color{white}\bfseries Printed result\\
-1 & true & throw then catch & -1 rejected\\
0 & false & Normal return & 0 accepted\\
100 & false & Normal return & 100 accepted\\
101 & true & throw then catch & 101 rejected\\
\end{tabularx}
\vspace{3mm}\footnotesize Follow the changing state in execution order.
\note{Manually derived trace for the complete solution. Write a validator that throws when marks are outside 0..100. Call the validator for each boundary example inside a try block. Print a success message only after validation returns normally.\par Syllabus reading: Deitel Ch 11. CLO-3.}
\end{frame}

\begin{frame}[fragile,t]{Solution result}
\begin{columns}[T,onlytextwidth]\begin{column}{0.60\textwidth}
-1 rejected\par 0 accepted\par 100 accepted\par 101 rejected
\end{column}\begin{column}{0.35\textwidth}\small
Why this result follows

Print a success message only after validation returns normally.
\end{column}\end{columns}
\note{Expected result: -1 rejected
0 accepted
100 accepted
101 rejected\par Syllabus reading: Deitel Ch 11. CLO-3.}
\end{frame}

\begin{frame}[fragile,t]{A common incorrect approach}
try \{ work(); \}\par catch (Exception ex) \{ \}\par catch (NumberFormatException ex) \{ \}
\vspace{1mm}\begin{block}{Example}
\begin{lstlisting}
The specific handler is unreachable. Catch NumberFormatException first, and give each handler a meaningful action.
\end{lstlisting}
\end{block}
\note{Ask students to predict the failure before discussing the explanation. The right side gives the correction.\par Syllabus reading: Deitel Ch 11. CLO-3.}
\end{frame}

\begin{frame}[fragile,t]{Boundary and failure checks}
\begin{columns}[T,onlytextwidth]\begin{column}{0.60\textwidth}
Check the stated input assumptions.

Read one integer and report malformed input using a specific catch. Keep parsing and range validation distinguishable.
\end{column}\begin{column}{0.35\textwidth}\small
Throw when marks\textless{}0 || marks\textgreater{}100. Inputs -1 and 101 fail. Inputs 0 and 100 return normally.
\end{column}\end{columns}
\note{Use the specification to derive each expected result. Exercise solution: Throw when marks\textless{}0 || marks\textgreater{}100. Inputs -1 and 101 fail. Inputs 0 and 100 return normally.\par Syllabus reading: Deitel Ch 11. CLO-3.}
\end{frame}

\begin{frame}[fragile,t]{Check your understanding}
\begin{columns}[T,onlytextwidth]\begin{column}{0.60\textwidth}
Does throws itself create an exception? Must unchecked exceptions be declared?
\end{column}\begin{column}{0.35\textwidth}\small
Write a prediction and explain the rule that supports it.
\end{column}\end{columns}
\note{Answer: No. throws declares a possible exception. Unchecked exceptions have no catch-or-declare compiler requirement.\par Syllabus reading: Deitel Ch 11. CLO-3.}
\end{frame}

\begin{frame}[fragile,t]{Answer and explanation}
\begin{columns}[T,onlytextwidth]\begin{column}{0.60\textwidth}
No. throws declares a possible exception. Unchecked exceptions have no catch-or-declare compiler requirement.
\end{column}\begin{column}{0.35\textwidth}\small
The specific handler is unreachable. Catch NumberFormatException first, and give each handler a meaningful action.
\end{column}\end{columns}
\note{Reconnect the answer to the relevant definition rather than treating it as a fact to memorise.\par Syllabus reading: Deitel Ch 11. CLO-3.}
\end{frame}


\begin{frame}[fragile,t]{Compare cases and predict outcomes}
\small\renewcommand{\arraystretch}{1.5}\rowcolors{2}{pale}{white}
\begin{tabularx}{\textwidth}{>{\raggedright\arraybackslash}X>{\raggedright\arraybackslash}X>{\raggedright\arraybackslash}X}\rowcolor{navy}
\color{white}\bfseries Token & \color{white}\bfseries Parsing outcome & \color{white}\bfseries Handler needed?\\
42 & int 42 & No\\
four & NumberFormatException & Yes\\
2147483648 & Outside int range & Yes\\
\end{tabularx}\par\vspace{3mm}\begin{exampleblock}{Discuss}Explain each expected result before running the example. Which case would reveal a wrong assumption?\end{exampleblock}
\note{Read the first two columns and invite predictions before discussing the outcome.}
\end{frame}

\begin{frame}[fragile,t]{Apply the idea}
\begin{alertblock}{Independent practice}Parse "four" inside try/catch, print "Invalid integer" on failure, then print "Finished" after the handler. Predict both output lines.\end{alertblock}\vspace{3mm}\begin{enumerate}\item Work through the input by hand.\item Write or correct the program.\item Justify the expected result.\end{enumerate}\note{Allow 3–5 minutes, then compare answers in pairs. Parse "four" inside try/catch, print "Invalid integer" on failure, then print "Finished" after the handler. Predict both output lines.}
\end{frame}

\begin{frame}[fragile,t]{Solution and reasoning}
\begin{lstlisting}
try {
  int n = Integer.parseInt("four");
  System.out.println(n);
} catch (NumberFormatException ex) {
  System.out.println("Invalid integer");
}
System.out.println("Finished");
\end{lstlisting}\vspace{1mm}\small\begin{itemize}
\item parseInt throws before n is printed.
\item The matching handler prints Invalid integer.
\item Normal control resumes after the try/catch and prints Finished.
\end{itemize}\note{Answer: parseInt throws before n is printed. The matching handler prints Invalid integer. Normal control resumes after the try/catch and prints Finished.}
\end{frame}

\begin{frame}[fragile,t]{Which exception is observed first?}
\begin{itemize}
\item One statement can contain more than one potential failure. Execution order determines the observed exception.
\item For a simple array assignment, the right{-}hand expression is evaluated before the final bounds check.
\item Only the first thrown exception reaches a matching handler in this execution.
\end{itemize}
\vspace{3mm}\footnotesize\textcolor{accent}{Source: supplied ISB campus slides. See speaker notes and ISB\_Source\_Map.md.}
\note{Adapted from the supplied ISB campus folder: \allowbreak Exception handling(updated).pptx, slides 29, 30. Made the divisor a variable and verified simple{-}assignment evaluation order against Java 17 JLS 15.26.1.}
\end{frame}

\begin{frame}[fragile,t]{Worked problem: Which exception is observed first?}
\begin{block}{Task}Trace int[] a=new int[5]; a[5]=30/zero with zero=0, then compare the case zero=2.\end{block}
\small Input: Values are fixed in the program for a reproducible demonstration.\par\vspace{3mm}\begin{exampleblock}{Before running}Predict the output and identify the variables that change.\end{exampleblock}
\note{Adapted from the supplied ISB campus folder: \allowbreak Exception handling(updated).pptx, slides 29, 30. Made the divisor a variable and verified simple{-}assignment evaluation order against Java 17 JLS 15.26.1.}
\end{frame}

\begin{frame}[fragile,t]{Complete ISB example (1/1)}
\begin{lstlisting}
public class IsbLecture24 {
  public static void main(String[] args) throws Exception {
    int[] a = new int[5];
    int zero = 0;
    try {
      a[5] = 30 / zero;
    } catch (ArithmeticException ex) {
      System.out.println("Arithmetic exception");
    } catch (ArrayIndexOutOfBoundsException ex) {
      System.out.println("Array index exception");
    }
    System.out.println("Rest of the code");
  }

}
\end{lstlisting}
\note{Adapted from the supplied ISB campus folder: \allowbreak Exception handling(updated).pptx, slides 29, 30. Made the divisor a variable and verified simple{-}assignment evaluation order against Java 17 JLS 15.26.1. Complete file: ISB\_Java\_Examples/IsbLecture24.java.}
\end{frame}

\begin{frame}[fragile,t]{Worked trace and expected result}
\small\renewcommand{\arraystretch}{1.4}\rowcolors{2}{pale}{white}
\begin{tabularx}{\textwidth}{>{\raggedright\arraybackslash}X>{\raggedright\arraybackslash}X>{\raggedright\arraybackslash}X}\rowcolor{navy}
\color{white}\bfseries Divisor / index & \color{white}\bfseries First failure & \color{white}\bfseries Handler\\
0 / 5 & Integer division by zero & Arithmetic\allowbreak Exception\\
2 / 5 & Invalid index after RHS & ArrayIndexOutOfBounds\allowbreak Exception\\
2 / 4 & None & Assignment succeeds\\
\end{tabularx}\par\vspace{2mm}
\begin{block}{Expected console output}\ttfamily\footnotesize Arithmetic exception\par Rest of the code\end{block}
\note{Adapted from the supplied ISB campus folder: \allowbreak Exception handling(updated).pptx, slides 29, 30. Made the divisor a variable and verified simple{-}assignment evaluation order against Java 17 JLS 15.26.1. Trace and expected values independently worked for this edition.}
\end{frame}

\begin{frame}[fragile,t]{Extend the ISB example}
\begin{alertblock}{Practice}What happens if catch(\allowbreak Exception ex) precedes both specific handlers?\end{alertblock}\vspace{3mm}\begin{exampleblock}{Explain your reasoning}Compilation fails because the later subclass handlers are unreachable.\end{exampleblock}
\note{Adapted from the supplied ISB campus folder: \allowbreak Exception handling(updated).pptx, slides 29, 30. Made the divisor a variable and verified simple{-}assignment evaluation order against Java 17 JLS 15.26.1. Answer: Compilation fails because the later subclass handlers are unreachable.}
\end{frame}
\begin{frame}[fragile,t]{Lecture summary}
\begin{columns}[T,onlytextwidth]\begin{column}{0.60\textwidth}
\begin{itemize}
\item propagation and the exception hierarchy
\item try, throw and catch
\item checked and unchecked exceptions
\end{itemize}
\end{column}\begin{column}{0.35\textwidth}\small
\begin{itemize}
\item Write a validator that throws when marks are outside 0..100.
\item Call the validator for each boundary example inside a try block.
\item Print a success message only after validation returns normally.
\end{itemize}
\end{column}\end{columns}
\note{Ask students to give one concrete example for the concepts listed. Use the worked solution to connect the topics.\par Syllabus reading: Deitel Ch 11. CLO-3.}
\end{frame}

\begin{frame}[fragile,t]{After-class practice}
\begin{columns}[T,onlytextwidth]\begin{column}{0.60\textwidth}
Read one integer and report malformed input using a specific catch. Keep parsing and range validation distinguishable.
\end{column}\begin{column}{0.35\textwidth}\small
Syllabus reading\par Deitel Ch 11

Next lecture: Cleanup and exception propagation
\end{column}\end{columns}
\note{Reading labels follow the supplied outline. Check topic headings against the available textbook edition. Require a program and expected results for the practice task.\par Syllabus reading: Deitel Ch 11. CLO-3.}
\end{frame}
\end{document}
